\begin{minted}{cpp}
int sa[1000020], pos[1000020], lcp[1000020];
int tmp[1000020], cnt[1000020], res[1000020];
void init(string s){
	s.push_back('#'); int sl = s.size();
	for (int i = 0; i < sl; i++){ sa[i] = i; }
	sort(sa, sa+sl, [&s](int s1, int s2){ return s[s1] < s[s2]; });
	res[sa[0]] = 0; for (int i = 1; i < sl; i++){
		int s1 = sa[i-1], s2 = sa[i]; res[s2] = res[s1] + (s[s1] != s[s2]);
	}
	for (int k = 1; ; k *= 2){
		for (int i = 0; i < sl; i++){ cnt[i] = 0; tmp[i] = (sa[i]-k%sl+sl)%sl; }
		for (int i = 0; i < sl; i++){ cnt[res[i]] += 1; }
		for (int i = 1; i < sl; i++){ cnt[i] += cnt[i-1]; }
		for (int i = sl-1; i >= 0; i--){
			int sp = tmp[i]; int r = res[sp];
			int p = cnt[r]; sa[p-1] = sp; cnt[r] -= 1;
		}
		tmp[sa[0]] = 0; for (int i = 1; i < sl; i++){
			int s1 = sa[i-1], s2 = sa[i];
			int t1 = (s1+k)%sl, t2 = (s2+k)%sl;
			tmp[s2] = tmp[s1] + (res[s1] != res[s2] || res[t1] != res[t2]);
		} memcpy(res, tmp, sizeof(res)); if (res[sa[sl-1]] == sl-1){ break; }
	} s.pop_back(); sl -= 1;
	for (int i = 0; i < sl; i++){ sa[i] = sa[i+1]; }
	for (int i = 0; i < sl; i++){ pos[sa[i]] = i; }
	int l = 0; for (int i = 0; i < sl; i++){
		l = max(l-1, 0); int p1 = pos[i], p2 = pos[i]-1; if (p2 < 0){ continue; }
		int s1 = sa[p1], s2 = sa[p2]; while (s1+l < sl && s2+l < sl){
			if (s[s1+l] != s[s2+l]){ break; } else{ l += 1; }
		} lcp[p2] = l;
	}
}
\end{minted}

\subsubsection*{Traversing}
\begin{minted}{cpp}
int idx[1000020]; // idx[i] = NULL position of sa[i]. Useful when S is multi-string.
// qry(st, ed) returns pair(min lcp[i], argmin) for [st, ed).
// solve on sa[st..ed][i..j).
ll solve(int st, int ed, int i){
	if (st == ed){ int j = idx[st]; return chk[st].fr <= ed && ed < chk[st].sc ? j-i : 0; }
	int j = qry(st, ed).fr; vector<int> mid{st};
	while (mid.back() < ed){
		pi2 p = qry(mid.back(), ed); if (p.fr != j){ break; }
		mid.push_back(p.sc+1);
	} mid.push_back(ed+1);
	ll ans = (chk[st].fr <= ed && ed < chk[st].sc ? j-i : 0);
	int len = mid.size(); for (int p = 1; p < len; p++){
		ans += solve(mid[p-1], mid[p]-1, j);
	} return ans;
}
\end{minted}